\section{Scope and limitations}
\label{introduction:scope-and-limitations}
% Academic rewrite
The study is bounded to vector workloads on Norwegian building polygons derived from conflated \acrshort{osm} and \acrshort{fkb} data, and excludes raster data and other feature types. It compares four configurations, DuckDB over GeoParquet, PostGIS, GeoPandas over Shapefile, and Apache Sedona on Databricks, on read performance alone, across point-in-polygon, k-nearest-neighbor, bounding-box, and spatial-join workloads; ingestion, update, and write performance are out of scope. Each workload is evaluated on run-time, network transfer, and operational cost across small, medium, and large dataset tiers, with the single-machine patterns run at the small and large tiers and the join across all three.

Several boundaries follow from holding the comparison to one controlled setting. The benchmark runs on a single cloud provider, Microsoft Azure, so the results do not establish how the configurations behave on other platforms. Operational cost reflects a pricing snapshot and the cost model assumption that intra-region, intra-tenant transfer is not billed as egress, and it shifts as vendor rates change. Queries are issued by a single client in sequence, leaving multi-user concurrency and contention unmeasured, and distributed scaling is observed only over 2 to 16 workers. No traditional format is tested in a distributed setting, since none targets one. These limitations bound how far the findings generalize; their effect on the validity of the measurements is treated in detail among the threats to validity in Section~\ref{research-design-and-methodology:threats-to-validity}.